% Worked Examples: Concept 1.3.4 - Linearisation
\documentclass[11pt,a4paper]{article}
\usepackage{worked-examples-template}

\title{\textbf{Worked Examples}\\[0.5em]
\large Concept 1.3.4: Linearisation\\[0.3em]
\normalsize \coursesubtitle{} -- \coursetitle}
\author{Assoc.\ Prof.\ William Robertson \and Dr Sean McGowan}
\date{\today}

\begin{document}

\maketitle
\tableofcontents
\newpage

\section{Concept 1.3.4: Linearisation}

\subsection{Example 1: Linearising a Simple Pendulum with Cart Using Small Angle Approximations}

This example demonstrates the complete linearisation process for a 3rd order nonlinear mechanical system using small angle approximations and Taylor series expansion. We'll work through the equations of motion, derive the nonlinear state equations, and systematically linearise them.

\paragraph{Problem:}
Consider a simple pendulum attached to a cart that can move horizontally. The system parameters are:
\begin{itemize}
\item Mass of cart: $M = 1.0$ kg
\item Mass of pendulum bob: $m = 0.3$ kg
\item Length of pendulum: $\ell = 0.5$ m
\item Gravitational acceleration: $g = 9.81$ m/s$^2$
\item Damping coefficient on cart: $c = 0.2$ N·s/m
\end{itemize}

The cart position is $q$ and the pendulum angle from vertical is $\theta$ (positive clockwise). A horizontal force $F$ is applied to the cart.

Starting from the equations of motion, derive:
\begin{enumerate}
\item The nonlinear state-space equations
\item The linearised state-space equations around the equilibrium point $q = 0$, $\theta = 0$ (pendulum hanging down)
\end{enumerate}

\paragraph{Solution:}

\textbf{Step 1: Derive equations of motion}

Using Lagrangian mechanics or force balance, the equations of motion are:
\begin{align}
(M + m)\ddot{q} + m\ell\cos\theta\,\ddot{\theta} - m\ell\sin\theta\,\dot{\theta}^2 + c\dot{q} &= F \label{eq:cart-eom}\\
m\ell^2\ddot{\theta} + m\ell\cos\theta\,\ddot{q} - m\ell g\sin\theta &= 0 \label{eq:pend-eom}
\end{align}

These are coupled second-order differential equations with nonlinearities in $\sin\theta$, $\cos\theta$, and the velocity coupling term $\dot{\theta}^2$.

\textbf{Step 2: Solve for accelerations to get explicit nonlinear ODEs}

We need to isolate $\ddot{q}$ and $\ddot{\theta}$. From equation~\eqref{eq:pend-eom}:
\begin{align}
\ddot{\theta} = -\frac{\cos\theta}{\ell}\ddot{q} + \frac{g}{\ell}\sin\theta
\end{align}

Substituting into equation~\eqref{eq:cart-eom}:
\begin{align}
(M + m)\ddot{q} + m\ell\cos\theta\left(-\frac{\cos\theta}{\ell}\ddot{q} + \frac{g}{\ell}\sin\theta\right) - m\ell\sin\theta\,\dot{\theta}^2 + c\dot{q} &= F\\
\left[(M + m) - m\cos^2\theta\right]\ddot{q} + mg\cos\theta\sin\theta - m\ell\sin\theta\,\dot{\theta}^2 + c\dot{q} &= F
\end{align}

Let $\Delta = M + m\sin^2\theta$. Then:
\begin{align}
\ddot{q} &= \frac{1}{\Delta}\left[F - c\dot{q} - mg\cos\theta\sin\theta + m\ell\sin\theta\,\dot{\theta}^2\right]\\
\ddot{\theta} &= \frac{1}{\ell\Delta}\left[(M+m)g\sin\theta - \cos\theta F + c\cos\theta\,\dot{q} - m\ell\cos\theta\sin\theta\,\dot{\theta}^2\right]
\end{align}

\textbf{Step 3: Define state variables}

Choose state vector $x = \begin{bmatrix} q & \theta & \dot{q} & \dot{\theta} \end{bmatrix}^T$. Note: This is actually a 4th order system, not 3rd order as stated in the problem. Let me proceed with the 4-state formulation, which is more typical for this system.

The nonlinear state equations are:
\begin{align}
\dot{x}_1 &= x_3\\
\dot{x}_2 &= x_4\\
\dot{x}_3 &= \frac{1}{M + m\sin^2 x_2}\left[F - c x_3 - mg\cos x_2\sin x_2 + m\ell\sin x_2\,x_4^2\right]\\
\dot{x}_4 &= \frac{1}{\ell(M+m\sin^2 x_2)}\left[(M+m)g\sin x_2 - \cos x_2 F + c\cos x_2\,x_3 - m\ell\cos x_2\sin x_2\,x_4^2\right]
\end{align}

\textbf{Step 4: Apply small angle approximations}

For the equilibrium point $x_e = \begin{bmatrix} 0 & 0 & 0 & 0 \end{bmatrix}^T$ and $F_e = 0$, we apply Taylor series approximations for small $\theta$:
\begin{align}
\sin\theta &\approx \theta - \frac{\theta^3}{6} + \cdots \approx \theta \quad \text{(first-order)}\\
\cos\theta &\approx 1 - \frac{\theta^2}{2} + \cdots \approx 1 \quad \text{(first-order)}\\
\sin\theta\cos\theta &\approx \theta \quad \text{(first-order)}\\
\sin^2\theta &\approx \theta^2 \approx 0 \quad \text{(second-order, neglected)}
\end{align}

Also note that for small deviations, the product terms like $\theta\dot{\theta}^2$ are third-order small and can be neglected in linear analysis.

\textbf{Step 5: Linearise the state equations}

Applying the approximations with $\theta \ll 1$, $\Delta \approx M + m(0) = M + m$:

For $\dot{x}_3$:
\begin{align}
\dot{x}_3 &\approx \frac{1}{M+m}\left[F - cx_3 - mgx_2 + 0\right]\\
&= -\frac{c}{M+m}x_3 - \frac{mg}{M+m}x_2 + \frac{1}{M+m}F
\end{align}

For $\dot{x}_4$:
\begin{align}
\dot{x}_4 &\approx \frac{1}{\ell(M+m)}\left[(M+m)gx_2 - F + cx_3 - 0\right]\\
&= \frac{g}{\ell}x_2 + \frac{c}{\ell(M+m)}x_3 - \frac{1}{\ell(M+m)}F
\end{align}

\textbf{Step 6: Write linearised state-space form}

The linearised system is $\dot{x} = Ax + Bu$ with input $u = F$:
\begin{align}
A &= \begin{bmatrix}
0 & 0 & 1 & 0\\
0 & 0 & 0 & 1\\
0 & -\dfrac{mg}{M+m} & -\dfrac{c}{M+m} & 0\\[0.5em]
0 & \dfrac{g}{\ell} & \dfrac{c}{\ell(M+m)} & 0
\end{bmatrix}, \quad
B = \begin{bmatrix}
0\\[0.2em]
0\\[0.2em]
\dfrac{1}{M+m}\\[0.5em]
-\dfrac{1}{\ell(M+m)}
\end{bmatrix}
\end{align}

\textbf{Step 7: Substitute numerical values}

With $M = 1.0$ kg, $m = 0.3$ kg, $\ell = 0.5$ m, $g = 9.81$ m/s$^2$, $c = 0.2$ N·s/m:
\begin{align}
A &= \begin{bmatrix}
0 & 0 & 1 & 0\\
0 & 0 & 0 & 1\\
0 & -2.265 & -0.154 & 0\\
0 & 19.62 & 0.308 & 0
\end{bmatrix}, \quad
B = \begin{bmatrix}
0\\
0\\
0.769\\
-1.538
\end{bmatrix}
\end{align}

\textbf{Key observations:}
\begin{itemize}
\item The linearisation process required: (1) deriving equations of motion, (2) explicit state form, (3) small angle approximations, (4) neglecting higher-order terms
\item The positive entry $A_{42} = g/\ell > 0$ indicates the system is unstable (pendulum falls away from vertical)
\item The linearisation is only valid near $\theta = 0$ (typically $|\theta| < 10°$ to $20°$)
\item All nonlinear coupling terms ($\dot{\theta}^2$, $\sin\theta\cos\theta$, etc.) are captured in linear form through Taylor expansion
\end{itemize}

\subsection{Example 2: Linearising a Magnetic Levitation System Using Jacobian Method}

This example demonstrates the systematic Jacobian approach to linearise a 4th order nonlinear system. The Jacobian method is more general and doesn't require manual Taylor series expansions.

\paragraph{Problem:}
Consider a magnetic levitation system where a ball of ferromagnetic material (mass $m = 0.05$ kg) is suspended by an electromagnet. The vertical position is $z$ (measured downward from the magnet, in meters) and the magnet current is $i$ (in Amperes).

The nonlinear dynamics are:
\begin{align}
m\ddot{z} &= mg - \frac{k_m i^2}{z^2}\\
L\dot{i} &= v - Ri
\end{align}
where:
\begin{itemize}
\item $k_m = 0.001$ N·m$^2$/A$^2$ (magnetic force constant)
\item $L = 0.1$ H (coil inductance)
\item $R = 1.0$ $\Omega$ (coil resistance)
\item $g = 9.81$ m/s$^2$ (gravitational acceleration)
\end{itemize}

The control input is the applied voltage $v$, and we can measure both position $z$ and current $i$.

Use the Jacobian method to linearise this system around an equilibrium point where the ball is suspended at $z_e = 0.01$ m.

\paragraph{Solution:}

\textbf{Step 1: Define state variables and write in standard form}

Define state vector $x = \begin{bmatrix} z & \dot{z} & i \end{bmatrix}^T = \begin{bmatrix} x_1 & x_2 & x_3 \end{bmatrix}^T$, input $u = v$, and output $y = \begin{bmatrix} z \\ i \end{bmatrix}^T$.

The nonlinear state equations in the form $\dot{x} = f(x,u)$ are:
\begin{align}
f(x,u) = \begin{bmatrix}
f_1(x,u)\\
f_2(x,u)\\
f_3(x,u)
\end{bmatrix} = \begin{bmatrix}
x_2\\[0.3em]
g - \dfrac{k_m x_3^2}{m x_1^2}\\[0.5em]
\dfrac{1}{L}(u - Rx_3)
\end{bmatrix}
\end{align}

The output equation in the form $y = h(x,u)$ is:
\begin{align}
h(x,u) = \begin{bmatrix}
h_1(x,u)\\
h_2(x,u)
\end{bmatrix} = \begin{bmatrix}
x_1\\
x_3
\end{bmatrix}
\end{align}

\textbf{Step 2: Find equilibrium point}

At equilibrium, all derivatives are zero: $\dot{x}_e = 0$. From $f_1 = x_2 = 0$, we have $\dot{z}_e = 0$.

From $f_2 = 0$:
\begin{align}
g - \frac{k_m i_e^2}{m z_e^2} &= 0\\
i_e^2 &= \frac{m g z_e^2}{k_m}\\
i_e &= \sqrt{\frac{m g z_e^2}{k_m}} = z_e\sqrt{\frac{mg}{k_m}}
\end{align}

With $z_e = 0.01$ m:
\begin{align}
i_e = 0.01\sqrt{\frac{0.05 \times 9.81}{0.001}} = 0.01\sqrt{490.5} = 0.2214 \text{ A}
\end{align}

From $f_3 = 0$:
\begin{align}
u_e = R i_e = 1.0 \times 0.2214 = 0.2214 \text{ V}
\end{align}

Therefore, the equilibrium point is:
\begin{align}
x_e = \begin{bmatrix} 0.01 \\ 0 \\ 0.2214 \end{bmatrix}, \quad u_e = 0.2214 \text{ V}
\end{align}

\textbf{Step 3: Compute Jacobian $A = \left.\frac{\partial f}{\partial x}\right|_{(x_e,u_e)}$}

The $A$ matrix is the Jacobian of $f$ with respect to $x$:
\begin{align}
A = \begin{bmatrix}
\dfrac{\partial f_1}{\partial x_1} & \dfrac{\partial f_1}{\partial x_2} & \dfrac{\partial f_1}{\partial x_3}\\[0.5em]
\dfrac{\partial f_2}{\partial x_1} & \dfrac{\partial f_2}{\partial x_2} & \dfrac{\partial f_2}{\partial x_3}\\[0.5em]
\dfrac{\partial f_3}{\partial x_1} & \dfrac{\partial f_3}{\partial x_2} & \dfrac{\partial f_3}{\partial x_3}
\end{bmatrix}
\end{align}

Calculate each partial derivative:

\emph{Row 1:} $f_1 = x_2$
\begin{align}
\frac{\partial f_1}{\partial x_1} = 0, \quad \frac{\partial f_1}{\partial x_2} = 1, \quad \frac{\partial f_1}{\partial x_3} = 0
\end{align}

\emph{Row 2:} $f_2 = g - \frac{k_m x_3^2}{m x_1^2}$
\begin{align}
\frac{\partial f_2}{\partial x_1} &= -\frac{k_m x_3^2}{m} \cdot \frac{\partial}{\partial x_1}\left(x_1^{-2}\right) = -\frac{k_m x_3^2}{m} \cdot (-2x_1^{-3}) = \frac{2k_m x_3^2}{m x_1^3}\\
\frac{\partial f_2}{\partial x_2} &= 0\\
\frac{\partial f_2}{\partial x_3} &= -\frac{k_m}{m x_1^2} \cdot 2x_3 = -\frac{2k_m x_3}{m x_1^2}
\end{align}

\emph{Row 3:} $f_3 = \frac{1}{L}(u - Rx_3)$
\begin{align}
\frac{\partial f_3}{\partial x_1} = 0, \quad \frac{\partial f_3}{\partial x_2} = 0, \quad \frac{\partial f_3}{\partial x_3} = -\frac{R}{L}
\end{align}

\textbf{Step 4: Evaluate at equilibrium point}

Substitute $(x_e, u_e)$:
\begin{align}
A_{21} &= \frac{2k_m i_e^2}{m z_e^3} = \frac{2 \times 0.001 \times (0.2214)^2}{0.05 \times (0.01)^3} = \frac{9.81 \times 10^{-8}}{5 \times 10^{-8}} = 1962 \text{ s}^{-2}\\
A_{23} &= -\frac{2k_m i_e}{m z_e^2} = -\frac{2 \times 0.001 \times 0.2214}{0.05 \times (0.01)^2} = -8.856 \text{ m/(A·s}^2\text{)}\\
A_{33} &= -\frac{R}{L} = -\frac{1.0}{0.1} = -10 \text{ s}^{-1}
\end{align}

Therefore:
\begin{align}
A = \begin{bmatrix}
0 & 1 & 0\\
1962 & 0 & -8.856\\
0 & 0 & -10
\end{bmatrix}
\end{align}

\textbf{Step 5: Compute Jacobian $B = \left.\frac{\partial f}{\partial u}\right|_{(x_e,u_e)}$}

\begin{align}
B = \begin{bmatrix}
\dfrac{\partial f_1}{\partial u}\\[0.3em]
\dfrac{\partial f_2}{\partial u}\\[0.3em]
\dfrac{\partial f_3}{\partial u}
\end{bmatrix} = \begin{bmatrix}
0\\[0.2em]
0\\[0.2em]
\dfrac{1}{L}
\end{bmatrix} = \begin{bmatrix}
0\\
0\\
10
\end{bmatrix}
\end{align}

\textbf{Step 6: Compute output matrices $C$ and $D$}

\begin{align}
C = \left.\frac{\partial h}{\partial x}\right|_{(x_e,u_e)} = \begin{bmatrix}
\dfrac{\partial h_1}{\partial x_1} & \dfrac{\partial h_1}{\partial x_2} & \dfrac{\partial h_1}{\partial x_3}\\[0.3em]
\dfrac{\partial h_2}{\partial x_1} & \dfrac{\partial h_2}{\partial x_2} & \dfrac{\partial h_2}{\partial x_3}
\end{bmatrix} = \begin{bmatrix}
1 & 0 & 0\\
0 & 0 & 1
\end{bmatrix}
\end{align}

\begin{align}
D = \left.\frac{\partial h}{\partial u}\right|_{(x_e,u_e)} = \begin{bmatrix}
0\\
0
\end{bmatrix}
\end{align}

\textbf{Step 7: Write the linearised system}

Define deviation variables:
\begin{align}
\delta x &= x - x_e, \quad \delta u = u - u_e, \quad \delta y = y - h(x_e,u_e)
\end{align}

The linearised system is:
\begin{align}
\delta\dot{x} &= A\delta x + B\delta u\\
\delta y &= C\delta x + D\delta u
\end{align}

Explicitly:
\begin{align}
\frac{d}{dt}\begin{bmatrix}
\delta z\\
\delta\dot{z}\\
\delta i
\end{bmatrix} &= \begin{bmatrix}
0 & 1 & 0\\
1962 & 0 & -8.856\\
0 & 0 & -10
\end{bmatrix}\begin{bmatrix}
\delta z\\
\delta\dot{z}\\
\delta i
\end{bmatrix} + \begin{bmatrix}
0\\
0\\
10
\end{bmatrix}\delta v\\
\begin{bmatrix}
\delta z\\
\delta i
\end{bmatrix} &= \begin{bmatrix}
1 & 0 & 0\\
0 & 0 & 1
\end{bmatrix}\begin{bmatrix}
\delta z\\
\delta\dot{z}\\
\delta i
\end{bmatrix}
\end{align}

\textbf{Step 8: Check stability}

The characteristic polynomial is:
\begin{align}
\det(sI - A) &= \det\begin{bmatrix}
s & -1 & 0\\
-1962 & s & 8.856\\
0 & 0 & s+10
\end{bmatrix}\\
&= (s+10)(s^2 - 1962)\\
&= (s+10)(s - 44.3)(s + 44.3)
\end{align}

Eigenvalues: $\lambda_1 = -10$, $\lambda_2 = 44.3$, $\lambda_3 = -44.3$.

Since $\lambda_2 > 0$, the open-loop system is \textbf{unstable}. This is expected: if the ball moves slightly down, gravity pulls it further down, and the magnetic force decreases rapidly (proportional to $1/z^2$).

\textbf{Key observations:}
\begin{itemize}
\item The Jacobian method is systematic: compute partial derivatives, evaluate at equilibrium
\item No manual small-angle approximations needed---the Jacobian automatically performs the linearisation
\item The large coefficient $A_{21} = 1962$ reflects the strong nonlinearity of the $1/z^2$ magnetic force
\item The system has one unstable pole requiring feedback control to stabilize
\item This approach extends easily to higher-order systems and more complex nonlinearities
\end{itemize}

\textbf{Verification note:} We can verify the equilibrium condition: at equilibrium, the magnetic force balances gravity:
\begin{align}
\frac{k_m i_e^2}{z_e^2} = mg \quad \checkmark
\end{align}
\end{document}
